\section{Классификация протоколов многосторонних конфиденциальных вычислений}
\label{sec:protocol-classification}

\subsection{Классификация по модели нарушителя}
\label{subsec:classification-threats}

\subsubsection{Получестная модель (semi-\allowbreak honest)}
В получестной модели (semi-honest; «честный, но любопытный» нарушитель) коррумпированный участник точно выполняет предписанные протоколом шаги, но анализирует полученные сообщения, пытаясь узнать дополнительные сведения о входах остальных участников. Коррумпированные участники могут объединить свои представления протокола: собственные входы и случайные значения, отправленные и полученные сообщения, а также результаты вычислений. В этой модели протокол считается безопасным, если такое представление не раскрывает о входах честных участников ничего сверх того, что следует из входов и выходов коррумпированных участников. Поскольку отклонения от протокола не рассматриваются, такие протоколы часто требуют меньше проверок и имеют меньшие накладные расходы, чем решения с защитой от активного нарушителя; точное сравнение зависит от протокола.

\subsubsection{Активная модель (malicious)}
В активной (злонамеренной) модели нарушитель может произвольно отклоняться от протокола: отправлять некорректные сообщения, менять ход вычисления или прекращать участие. Протокол с защитой от такого противника должен сохранять конфиденциальность и гарантировать, что выданный честным сторонам результат соответствует вычисляемой функции на допустимых входах; обнаруженное нарушение может привести к прерыванию вычислений. Однако активная модель безопасности сама по себе не гарантирует справедливость: нарушитель может узнать свой результат и затем прервать протокол, не позволив честным участникам получить результаты \cite{06}. Набор гарантий зависит от протокола, допустимого числа сговорившихся участников и сетевых допущений. Для контроля отклонений применяют аутентифицированные доли с MAC, доказательства с нулевым разглашением и технику cut-and-choose для искаженных схем; эти механизмы увеличивают вычислительные и коммуникационные затраты.

\subsubsection{Скрытная модель (covert)}
Скрытная (covert) модель занимает промежуточное положение между получестной и активной: нарушитель может отклоняться от протокола, но попытка обмана, нарушающая безопасность, обнаруживается с вероятностью не ниже заданного параметра \(\varepsilon\). Предполагается, что риск разоблачения и возможные санкции сдерживают злоумышленника. Значение \(\varepsilon\) выбирают при проектировании, универсального диапазона нет. Модель может быть уместна в коммерческих сценариях с репутационными, юридическими или финансовыми последствиями и позволяет снизить издержки по сравнению с полной активной защитой \cite{06}.

\subsection{Классификация по количеству участников}

\subsubsection{Двусторонние протоколы (2PC)}
Двусторонний вариант МКВ предполагает, что только два участника \(P_A\) и \(P_B\) совместно вычисляют функцию \(f(x_A,x_B)\).  Для двух сторон применимы техники, не масштабируемые на множество участников.  Используется схема зашифрованных схем Яо, в которой функция описывается в виде булевой схемы, затем «garbler» отправляет «вычислителю» зашифрованную схему, после чего обе стороны получают результат, ничего более не узнавая.  Яо-протокол обеспечивает безопасность против полудобросовестного противника и характеризуется постоянным числом раундов общения, независимо от сложности функции.  Расширения этого подхода включают активную защиту с cut-and-choose и протокол BMR, который позволяет использовать garbled circuits в многостороннем случае.

\subsubsection{Многосторонние протоколы (n-party)}
Когда вычисление выполняют более двух участников, роли garbler и evaluator исчезают; данные каждой проводки схемы представляются в виде долей, распределённых между сторонами.  В таких протоколах функция изображается как арифметическая схема над конечным полем, а операции сложения и умножения выполняются над секретными долями.  Типичный пример --- протокол GMW, разработанный Голдрайхом, Микали и Вигдерсоном; он использует аддитивное разделение секрета и масштабируется на произвольное число сторон.  GMW допускает активных нарушителей, если большинство участников честное, и вводит доказательства с нулевым разглашением для защиты от обманов.  Другой пример --- протокол BGW, который строится на схеме Шамира и использует арифметические схемы.  Он эффективен для функций, описываемых операциями сложения и умножения, и позволяет устанавливать пороговое значение количества нечестных участников.

\subsubsection{Допущение о честном большинстве}
В протоколах на основе разделения секрета важную роль играет порог \(t\) числа нечестных сторон.  Для схемы Шамира секрета величина порога, при котором обеспечивается информационно-\allowbreak теоретическая конфиденциальность, зависит от модели противника: при пассивном нарушителе условие \(t< n/2\) достаточно для восстановления секрета, а при активном нарушителе требуется \(t< n/3\).  Аддитивное разделение секрета может выдерживать до \(n-1\) нарушителей в обеих моделях, но требует дополнительных механизмов для проверки корректности (например, MAC-аутентификации в протоколе SPDZ).  Таким образом, классификация по количеству участников тесно связана с предположением об относительном числе честных сторон, что влияет на выбор схемы и гарантии безопасности.

\subsection{Классификация по используемому методу}
\label{subsec:classification-primitives}

В зависимости от используемых криптографических примитивов протоколы МКВ подразделяют на несколько групп. Механизмы и сравнительная характеристика примитивов приведены в разделе~\ref{sec:cryptographic-foundations}. Критерии согласуются с обзором П. Морейры, сопоставляющим протоколы по предположениям безопасности, эффективности, удобству использования и функциональным возможностям \cite{10}.

Протоколы на основе разделения секрета используют секретные доли для представления данных. К этой группе относятся BGW, GMW и SPDZ. Гарантии безопасности каждого протокола зависят от принятой модели противника.

К протоколам на основе зашифрованных схем относятся протокол Яо и его многосторонние расширения, например BMR. Отдельную группу составляют протоколы на основе гомоморфного шифрования, в которых данные обрабатываются посредством операций над шифротекстами.

Гибридные протоколы сочетают несколько криптографических примитивов на разных этапах вычисления. Классификация по основному методу не исключает вспомогательных механизмов, в частности ослеплённой передачи (см. раздел~\ref{sec:cryptographic-foundations}).

Специализированные протоколы для приватного пересечения множеств, аукционов и других задач выделяют по назначению, а не по используемому криптографическому методу. Они могут относиться к любой из перечисленных групп. Обзор PSI-протоколов Д. Моралеса, И. Агудо и Х. Лопеса выделяет гомоморфное шифрование, ослеплённую передачу и забывающие псевдослучайные функции и сопоставляет решения по сложности и практической применимости \cite{09}. Классификация по назначению дополняет, но не заменяет критерии примитивов, числа участников и модели противника.

\subsection{Вывод}

Критерии выбора протокола заданы классификацией (см. подразделы~\ref{subsec:classification-threats}--\ref{subsec:classification-primitives}). Двусторонние схемы (garbled circuits, FHE) обеспечивают низкую латентность при сравнении и обработке конфиденциальных данных двух сторон; многосторонние GMW, BGW и SPDZ применяются для распределённых вычислений при честном большинстве. Зависимость коммуникационных и вычислительных затрат от примитивов отражена в таблице~\ref{tab:mpc-primitives}. В следующем разделе рассмотрены применения МКВ в аналитике данных, МО, финансах и медицине.